\begin{table}[H]
\centering
\caption*{\textbf{Quadro A.1 - Contratos de catálogo, plano e routing}}
\scriptsize
\renewcommand{\arraystretch}{1.1}
\begin{tabularx}{\linewidth}{|>{\raggedright\arraybackslash}p{0.17\linewidth}|>{\raggedright\arraybackslash}p{0.34\linewidth}|Y|}
\hline
\textbf{Objeto} & \textbf{Campos mínimos} & \textbf{Invariantes} \\
\hline
\texttt{SkillRecord} &
\texttt{name}, \texttt{description}, \texttt{body},
\texttt{allowedTools[]}, \texttt{indexText} &
nome único; body não vazio; todas as tools resolvem no registro \\
\hline
\texttt{ToolDescriptor} &
\texttt{name}, \texttt{inputSchema}, \texttt{outputSchema} &
nome único; schemas disponíveis ao executor \\
\hline
\texttt{GoalNode} &
\texttt{id}, \texttt{goal}, \texttt{doneWhen[]}, \texttt{dependsOn[]},
\texttt{status}, \texttt{deliver} &
ID único; dependências válidas; status permitido; objetivo orientado a resultado \\
\hline
\texttt{Plan} & \texttt{nodes[]}, \texttt{revision} &
DAG acíclico; revisão inteira não negativa, contígua e atribuída pelo runtime;
ao menos um nó entregável; nós concluídos imutáveis em revisões \\
\hline
\texttt{SkillHint} &
\texttt{goalId}, \texttt{skillName}, \texttt{evidence}, \texttt{effect} &
fundamentada no body; \texttt{effect} indica vocabulário, lacuna, divisão ou dependência \\
\hline
\texttt{SkillCandidate} &
\texttt{skillName}, \texttt{score}, \texttt{rank}, \texttt{rationale} &
aponta para skill canônica; rank total e determinístico após desempate \\
\hline
\texttt{OrderedBundle} &
\texttt{goalId}, \texttt{skills[]}, \texttt{selectionRationale} &
cardinalidade de zero a $K_{\max}$; sem duplicatas; ordem final observável \\
\hline
\end{tabularx}
\par\smallskip\raggedright\small Fonte: elaboração própria (2026).
\end{table}
